\scanpage{85}
où $p_{X'}$ désigne le morphisme structural $X'\quad X$. D'autre part, pour donner un sens à 4.1. (ii) et justifier \deleted{\unsure{sa|la}}la démonstration indiquée de 4.1. , il y a lieu d'intorduire un univers \quad tel que \qquad \unsure{,|;} et considérer les faisceaux intervenant dans les formules \deleted{\unsure{5}}4.1. (i\unsure{i)} comme des faisceaux à valeurs dans \quad -(Ens).

\uline{Proposition} 4L1. (i) Le foncteur $u^x$ est pleinement fidèle, donc pour tout faisceau $G$ sur $X_{\text{ét}}$, l'homomorphisme canonique $G\qquad u_x u^x(G)$ est un isomorphisme.

(ii) Soient $G$ (resp. $F$) un faisceau abélien sur $X_{\text{ét}}$ (resp. $(\mathrm{Sch})_{/X}$). Alors les homomorphismes canoniques (définis par exemple comme edge-homomorphismes de la suite \unsure{spectrahe|spectrale} de Leray V 5.5.1.) suivants sont\deleted{\unsure{s|x}} des isomorphismes:

\[
\begin{aligned}
H^{\deleted{\unsure{x|{}}}*}(X_{\text{ét}},F\,u)&\simeq H^{\deleted{x}*}((\mathrm{Sch})_{/X},F)\\
H^*(X_{\text{ét}},G)&\simeq H^{\deleted{x}*}((\mathrm{Sch})_{/X},u^{\deleted{x}*}(G))\ .
\end{aligned}
\]

En effet, le foncteur $u$ est un comorphisme (VI 7.3.6.) et on peut alors appliquer VI 7.3.7. et VI 7.3.9.
% archive sheet number: 84

